\section{Tensor conventions and occupation}\label{sec:preliminaries}

Let $V=\Mat$, equipped with the nondegenerate coordinate pairing
\[
 \langle A,X\rangle=\sum_{i,j=0}^2 A_{ij}X_{ij}=\tr(A^TX).
\]
All matrix and vector coordinates are indexed by $0,1,2$. Write
$e_0=(1,0,0)^T$, $e_1=(0,1,0)^T$, $e_2=(0,0,1)^T$ and
$E_{ij}=e_i e_j^T$. Throughout the paper,
\begin{equation}\label{eq:tensor}
 \T=\sum_{i,j,k=0}^2 E_{ij}\otimes E_{jk}\otimes E_{ik}
   \in V\otimes V\otimes V.
\end{equation}
Thus its coefficient at $(3i+j,3j+k,3i+k)$ is $1$, and all other
coefficients are zero. Its contraction against $X,Y,Z$ is
$\tr(XYZ^T)$.

\begin{definition}
A decomposition of length $r$ is an identity
\begin{equation}\label{eq:decomp}
 \T=\sum_{t=1}^r A_t\otimes B_t\otimes C_t,
 \qquad A_t,B_t,C_t\in V.
\end{equation}
The tensor rank $R_{\F}(\T)$ is the least such $r$. Unless explicitly
excluded, zero summands and repeated factors are permitted in a
length-$r$ expression. A minimal expression has no zero summands.
\end{definition}

By contraction of~\eqref{eq:decomp}, the corresponding algorithm is
\begin{equation}\label{eq:algorithm}
 XY=\sum_{t=1}^r\langle A_t,X\rangle
                      \langle B_t,Y\rangle C_t
 \qquad(X,Y\in V).
\end{equation}
Conversely, evaluating~\eqref{eq:algorithm} on pairs of matrix units
recovers every coefficient of~\eqref{eq:decomp}. Matrix rank
$\rank(A_t)$ and tensor rank $R_{\F}(\T)$ will always be distinguished.

For a linear subspace $W\leq V$, put
\[
 q_W:V\longrightarrow V/W,\qquad
 T_W=(q_W\otimes\id\otimes\id)\T.
\]
The dual of $V/W$ is naturally $W^\perp$ under the coordinate pairing.
Consequently $T_W$ represents multiplication with its first input
restricted to $W^\perp$, not to $W$. Indeed, a linear functional on
$V/W$ pulls back to a functional vanishing on $W$, and every such
functional factors uniquely through $q_W$. Applying this identification
in the first tensor slot proves the assertion, including equality of
the two tensor ranks.

\begin{lemma}[Occupation inequality]\label{lem:occupation}
For a length-$r$ decomposition, let
$m(W)=|\{t:A_t\in W\}|$. Then
\begin{equation}\label{eq:occupation}
 m(W)\leq r-R_{\F}(T_W).
\end{equation}
More generally, if $W\leq U$ and $T_W$ has a length-$r$
decomposition, at most $r-R_{\F}(T_U)$ first factors, counted with
multiplicity in $V/W$, belong to $U/W$.
\end{lemma}
\begin{proof}
Applying $q_W$ to~\eqref{eq:decomp} kills every indexed term
with first factor in $W$. Discarding these terms gives a decomposition
of $T_W$ with at most $r-m(W)$ terms. For the
second assertion apply the induced map $V/W\to V/U$ instead.
\end{proof}

\subsection{Tensor symmetries}
For $U,V_0,H\in\GL_3(\F)$ the coefficient transformation
\begin{equation}\label{eq:symmetry}
 (A,B,C)\longmapsto
 (U^TAV_0^{-T},\ V_0^TBH^{-T},\ U^{-1}CH)
\end{equation}
preserves~\eqref{eq:tensor}. To check the contragredient factors,
consider the input transformation
\[
 (X,Y,Z)\longmapsto
 (UXV_0^{-1},\ V_0YH^{-1},\ U^{-T}ZH^T).
\]
It preserves $\tr(XYZ^T)$ by cancellation and cyclicity of trace.
Pullback under the coordinate pairing gives~\eqref{eq:symmetry}.
The transformation
\begin{equation}\label{eq:transpose}
 (A,B,C)\longmapsto(A^T,C,B)
\end{equation}
is also a symmetry, as follows directly by exchanging $i$ and $j$ in
the transformed sum~\eqref{eq:tensor}. These symmetries preserve the
matrix ranks of the factors and quotient tensor ranks.

The induced group on $V$ consists of $A\mapsto PAQ$ and
$A\mapsto PA^TQ$, for invertible $P,Q$. We denote it by $G$.
It has $2\cdot168^2=56{,}448$ elements. Its action on subspaces is the
one used in the frozen catalogue. Each of these first-factor actions
extends to a symmetry of the complete tensor by~\eqref{eq:symmetry}
and~\eqref{eq:transpose}.
